% chapters/06_grenzen.tex

\chapter{Grenzen}
\label{chap:grenzen}
\label{chap:diskussion}
\label{sec:limitierungen}
\label{sec:validitaet}

\paragraph{Stichprobe.} Closed-Loop-Messungen umfassen 10, 20 oder 60 Episoden. Für 0/20 reicht
das 95-\%-Intervall bis 17\,\%, für 2/60 von 0,4 bis 11,5\,\%. Griffraten zwischen 0 und 2 von
20 sind deshalb nicht voneinander zu trennen; belastbar sind nur große Unterschiede wie bei
\texttt{touched} und der Fingerspanne.

\paragraph{Simulation ist nicht der Roboter.} Alle Aussagen zur Aufgabenerfüllung stammen aus
der Simulation. Die Ursache des 3-cm-Versatzes ist unbekannt und nur empirisch korrigiert; die
Korrektur ist ein Mittelwert: Das Orakel hebt in 5 der 7 Lagen 240--246, danach nur in 1 von
13 (Lauf~69). Ob der Verlust von N1.6 auf Echtbildern nach dem Co-Training
(\cref{sec:erg-cotrain}) am Roboter spürbar ist, kann nur ein Robotertest zeigen.

\paragraph{Grenzen der Metriken.}
\label{sec:grenzen-metriken}
\label{sec:messartefakte}
\begin{itemize}
  \item \texttt{lifted} zählt auch einen gekippten Würfel (Schwerpunkt an der Kante +1,04~cm,
    an der Ecke +1,83~cm). Deshalb gibt es seit dem 13.10.\ \texttt{lifted\_clear}.
  \item \texttt{lifted\_clear} zählt auch einen Würfel, der auf einem anderen liegt, während
    eine Hand nah ist (Lauf~94). Jeder Treffer braucht den Videobefund.
  \item Der SigLIP-Abstand~\cite{zhai2023siglip} (\texttt{gap}) misst nur die Szene, nicht den Encoder der Policy;
    als Kontrolle für das Co-Training ist er ungeeignet (Lauf~96).
  \item Die 80-\%-Schwelle des visuellen Gates lassen auch alte Modelle ohne Co-Training hinter
    sich (Lauf~87); eine Schwelle relativ zum besten nicht co-trainierten Modell derselben
    Generation wäre strenger.
  \item Bis zum 30.10.\ war die Turmhöhe für \texttt{stacked} fest auf 8~cm gesetzt. Mit
    4-cm-Würfeln (Läufe 88, 93, 97) misst ein perfekter Dreierturm genau 8~cm und zählte damit
    kaum. Das ändert die Aussage \enquote{nie gestapelt} nicht: In keiner Episode wurde ein
    Würfel auf einen anderen gesetzt. Seitdem gilt zwei Kanten minus 2~cm.
  \item Kennzahlen, die ein Maximum über die Episode bilden (Fingerspanne), hängen von der
    Episodenlänge ab. Verglichen wird nur bei gleicher Länge.
\end{itemize}

\paragraph{Ablauf der Messungen.} Closed Loops hingen dreimal ohne Fehlermeldung mitten in einer
Episode (Läufe 89, 90, 97); unvollständige Versuche wurden nicht gewertet. Lauf~85 lief
versehentlich ohne \texttt{SIM\_CUBE\_EDGE}. Auf dem Sim-Server lag der Datensatz zeitweise
gemischt aus v2.1- und v3.0-Dateien vor (Lauf~56); welche früheren Läufe davon betroffen
waren, ist nicht geklärt.

\paragraph{Vergleich N1.6/N1.7.} N1.7 lief mit dem NVIDIA-Rezept (State-Dropout 0,2, andere
Normierung). Unterschiede zwischen den Generationen sind deshalb nicht allein dem Modell
zuzuschreiben.
